\subsection{相平衡条件与稳定相的选择}
\label{th:phasestability}
同温同压的两相处于机械、热学平衡后，仍可能发生物质转移。设从相1向相2可逆转移 $\dd\nu$ 摩尔，总 Gibbs 自由能变化为
\[
 \dd G=(g_2-g_1)\dd\nu.
\]
若 $g_2<g_1$，向相2转移会降低总 $G$，因此相2更稳定；只有 $g_2=g_1$ 时，两相才可在不产生净转移趋势的条件下共存。因而相平衡除了同温、同压，还要求转移物质时两相之间没有自由能降低的方向。

沿共存线改变温度、压强时，两相的 $g$ 继续相等。前文 Clapeyron 方程就是这一相等条件的全微分。斜率的分子是相变熵差，分母是体积差：熵差规定升温倾向哪一相，体积差规定加压倾向哪一相。对通常的液气相变，升温有利于气相，加压有利于液相，两种影响共同形成向右上升的饱和曲线。

\Diagram{fig:phasegibbs}{两相摩尔自由能在共存压强处相等}{
 \draw[axis](0,0)--(5.5,0)node[right]{$p$};
 \draw[axis](0,0)--(0,3.5)node[above]{$g$};
 \draw[curve](.6,1.5)--(4.8,2.2)node[right]{$g_l$};
 \draw[curve](.6,.65)--(4.8,2.95)node[right]{$g_g$};
 \draw[guide](2.831,0)--(2.831,1.872);
 \node[below]at(2.831,0){$p_s(T)$};
 \node at(2.2,3.1){固定温度，斜率为摩尔体积};
}
图\ref{fig:phasegibbs}示意共存点附近的变化。因 $g_p|_T=v$，气相的曲线斜率远大于液相；低于饱和压强时气相自由能较低，高于饱和压强时液相较低。图中的直线只表示局部线性化，跨越较大压强区间时应按各相体积积分。

\subsection{两相平台、杠杆规则与可压缩性}
固定温度下，若每摩尔系统的平均体积 $v$ 位于 $v_l$ 与 $v_g$ 之间，稳定平衡可由液气两相组成。令气相摩尔分数为 $x$，则
\[
 v=(1-x)v_l+xv_g,\qquad
 x=\frac{v-v_l}{v_g-v_l}.
\]
压强保持 $p_s(T)$，体积增大通过增加气相比例实现。此时形式上的等温压缩率很大，反映相比例调整，不能再用单一均匀相的微小密度响应解释。若相变成核受阻，实际系统可能暂时停留在亚稳态，但这需要额外的动力学条件。

\Diagram{fig:leverrule}{共存平台上的总体积由两端相体积加权}{
 \draw[axis](0,0)--(6,0)node[right]{$v$};
 \foreach \x/\t in{1/v_l,3.3/v,5/v_g}{
  \draw(\x,-.13)--(\x,.13);\node[below]at(\x,-.15){$\t$};
 }
 \draw[<->](1,.7)--(3.3,.7)node[midway,above]{$x(v_g-v_l)$};
 \draw[<->](3.3,1.5)--(5,1.5)node[midway,above]{$(1-x)(v_g-v_l)$};
}
图\ref{fig:leverrule}中，靠近液相端点意味着气相比例小。本式使用摩尔体积，权重相应为摩尔分数；采用质量分数或体积分数时，应连同被加权的量一起变换。计算两相热容时，显热项与 $x(T)$ 的变化项都从这一约束出发。

\subsection{等面积条件与自由能相等}
在可解析延拓的状态方程中，沿一条等温线有 $\dd g=v\dd p$。假设两端 $v_l,v_g$ 对应同一压强 $p_s$，积分并分部积分，
\[
 g_g-g_l=\int_{v_l}^{v_g}v\,p'(v)\dd v
 =p_s(v_g-v_l)-\int_{v_l}^{v_g}p(v)\dd v.
\]
要求两端自由能相等，便得到
\[
 \int_{v_l}^{v_g}[p(v)-p_s]\dd v=0.
\]
这正是等面积条件。积分中使用的中间均匀态分支是状态方程的数学延拓，不表示系统在平衡相变时真实经过它们。真实等温压缩沿共存平台改变相比例；曲线在平台上下的有向面积相抵，是两端化学势相等的另一种表达。

\subsection{饱和蒸气压近似的使用层次}
从 Clapeyron 方程到常见的指数蒸气压式，需要依次采用三个近似：液相比气相体积可忽略；气相满足理想气体状态方程；潜热在所讨论温区内近似不变。前两项把斜率化为 $L_m p_s/(RT^2)$，最后一项才允许把 $L_m$ 提到积分号外。

温度接近临界点时，两相体积差和潜热同时趋小，忽略液相体积的近似失效；在高压区，气相也可能明显偏离理想气体。对较宽温区，即使气相足够稀薄，也应保留 $L_m(T)$ 的变化。后面的变潜热模型只改进第三层近似，仍保留前两层，所得幂律乘指数式也应在这一范围内理解。

